\chapter{Abroad}
\label{ch:internationalism}

\section{Why a small poor island had a foreign policy}

For thirty years Cuba conducted a foreign policy several orders of magnitude
larger than its size, and any account of the period that treats this as
peripheral has missed how much of the country's resources, self-understanding
and reputation went into it. At its peak Cuba had tens of thousands of soldiers
in Africa, doctors and teachers in dozens of countries, and a claim on the
attention of the non-aligned world that no state of six million people has had
before or since.

Three motives ran together and it is not always possible to separate them.
There was conviction: the leadership genuinely believed in a global revolution
and acted on that belief at real cost. There was security: a country under
permanent threat from a superpower ninety miles away buys allies and buys
standing, and standing is a form of deterrence. And there was the patron: after
1972 Cuban military deployments in Africa were logistically impossible without
Soviet transport and finance, and they served Soviet purposes, which is not the
same as saying Moscow ordered them --- the Angolan intervention in particular
appears to have been decided in Havana and presented to Moscow afterwards.

\section{The guerrilla decade}

Through the 1960s Cuba supported armed insurgency across Latin America, in
Venezuela, Guatemala, Colombia, Peru, Bolivia and Nicaragua, and in Africa.
Almost all of it failed. The doctrine, elaborated by R\'egis Debray from
Guevara's writings, held that a small armed nucleus in the countryside could
create the conditions for revolution rather than waiting for them, which was a
misreading of Cuba's own case: the Sierra Maestra had succeeded alongside an
urban movement, a collapsing army and an isolated dictator, none of which the
doctrine required.

Guevara left Cuba in 1965 to apply it, first in the Congo, where the campaign
was a failure he documented himself with striking honesty, and then in Bolivia,
where he was captured and killed by the Bolivian army with CIA assistance in
October 1967. His death converted him into the most reproduced image of the
twentieth century and removed from Cuban politics the one figure whose prestige
was independent of Castro's.

The policy also cost Cuba with Moscow, which was pursuing peaceful coexistence
and had working relations with several of the governments Cuba was trying to
overthrow, and with the Latin American communist parties, which Havana was
denouncing as passive. The 1966 Tricontinental Conference in Havana was the
high point of Cuban third-worldism and close to the low point of Cuban--Soviet
relations.

The alignment was purchased back in 1968. In August, when Warsaw Pact tanks
entered Czechoslovakia, Castro went on television and endorsed the invasion. He
did so in a speech that conceded almost everything --- that the invasion
violated legality, that the socialist camp had no right to do it under the
rules it claimed --- and endorsed it anyway on the grounds that the alternative
was Czechoslovakia's drift to the West. The speech is the most revealing
document of Cuban foreign policy in the period, and the trade it embodied was
explicit within months: Soviet oil deliveries were increased, the economic
integration of the 1970s followed, and Cuban support for guerrilla movements in
Latin America was wound down.

\section{Africa}

Cuba's military engagement in Africa began in Algeria in 1963 and became large
in Angola in 1975.

When Portugal's empire collapsed, Angola's independence was contested by three
movements. The Soviet-aligned MPLA held Luanda; the FNLA advanced from Zaire;
and UNITA was backed from the south by South Africa, which invaded in October
1975 with an armoured column heading for the capital. Cuba's response,
Operation Carlota, began in November \citep{gleijeses2002}: troops flown in on Cuban and later Soviet
aircraft, stopping the South African advance short of Luanda. Within months
there were 36,000 Cuban soldiers in Angola. They stayed for sixteen years. Over
the whole period something like 300,000 Cuban military personnel and 50,000
civilians served there \citep{gleijeses2013}.\unv{}

The war's decisive phase came in 1987--88 around Cuito Cuanavale in the
south-east, where a Cuban and Angolan force stopped a South African offensive,
and then in a Cuban armoured push toward the Namibian border that changed the
military balance in the south-west. The battle's outcome is disputed in detail
--- both sides claim it --- but its consequence is not: the quadripartite
negotiations that followed produced South African withdrawal from Angola,
Namibian independence in 1990, and Cuban withdrawal, and they took place
because Pretoria concluded it could no longer hold the line militarily. Nelson
Mandela, released in 1990, went to Havana in 1991 and said so directly. Whatever
else is in the ledger, Cuba's contribution to ending South African control of
Namibia and to the pressure that ended apartheid is real, was paid for in Cuban
lives, and was not done for money.

Ethiopia is the other half of the record and it is much worse. In 1977--78
Cuba sent some 16,000 troops to defend Ethiopia against a Somali invasion of
the Ogaden --- defensible on its own terms, since Somalia was the aggressor and
Cuba refused to join the subsequent offensive into Somalia proper. But the
government being defended was the Derg of Mengistu Haile Mariam, which was in
the middle of the Red Terror, and Cuba's presence continued through it and
through the famine and the forced resettlements of the 1980s. Cuba had also
been the ally of Somalia and of the Eritrean independence movement before
switching sides, and the Eritreans, who had trained in Cuba, were now on the
other end of the guns. Whatever internationalism meant in Angola, in Ethiopia
it meant supporting a mass-murdering client of the same patron.

The official Cuban count of dead across all internationalist missions,
published in 1991, was 2,289.\unv{} The families were told what was necessary
and the coffins came home together, in a single ceremony, sixteen years of them
at once, in December 1989.

\subsection{The doctors}

Running in parallel, and outliving all of it, was medical internationalism.
The first Cuban medical brigade went to Algeria in 1963, fifty-six people, sent
by a country that had just lost half its own physicians \citep{kirk2015}. By the 1980s Cuban
doctors were working across Africa, Central America and the Caribbean, mostly
without charge, in places where no other country would send anyone. The Latin
American School of Medicine, founded later, trained foreign students free.

This is the origin of the reputation that Chapter~\ref{ch:life} will argue is
Cuba's most valuable intangible asset --- and, from the 2000s, of the country's
largest export earner, which is a different thing and carries its own problems
of pay, consent and conditions that Part II will take up honestly.

\section{Mariel}

In April 1980 a bus crashed through the fence of the Peruvian embassy in
Havana. The guards were withdrawn after a dispute over asylum, and within
seventy-two hours around 10,800 people were inside the embassy grounds. The
scale of it --- ten thousand people in a garden, in a country whose government
had said for twenty years that only a handful of \emph{gusanos} wanted to leave
--- was a shock from which the official narrative never quite recovered.

Castro's response was to open the port of Mariel and invite Cuban Americans to
come and collect their relatives by boat. Between April and October 1980 about
125,000 people crossed to Florida. The Cuban state used the opening to empty
some of its prisons and psychiatric hospitals into the boats --- the number is
disputed but the practice is not, and it was deliberate, both to rid itself of
a burden and to discredit the emigration as criminal. It also organised
\emph{actos de repudio}, mob demonstrations outside the homes of those leaving,
in which neighbours and co-workers were assembled to shout, throw eggs and
sometimes beat people who had announced their departure. Cubans were made to
brutalise their neighbours as a condition of their own standing.

The consequences ran both ways. In the United States, Mariel arrived in an
election year, overwhelmed Florida's capacity, was fixed in the public mind by
the criminal contingent, and did lasting damage to the standing of Cuban
immigrants generally. In Cuba, the memory of the repudiation mobs is one of the
sharpest wounds the revolution inflicted on its own society --- not because of
what the state did, but because of what it made ordinary people do to each
other, which is harder to forgive and harder to forget. Any reconciliation
process proposed in Part III has to have an answer for the \emph{acto de
repudio}, and the answer is not that it was long ago.

\section{The end of the decade}

Two events closed the period.

In June and July 1989 General Arnaldo Ochoa, a hero of the Angolan campaign and
one of the three or four most prominent officers in the country, was arrested,
tried and executed along with three others, and a group of senior Interior
Ministry officials including the de la Guardia brothers was convicted, on
charges of drug trafficking and corruption. The Interior Ministry was purged and
placed under military control. That the trafficking occurred is not seriously
disputed. Whether the trial was about trafficking is: Ochoa was popular, was
associated with the officers who had seen how the Soviet Union was changing, and
had reportedly been critical of the leadership. The proceedings were televised,
the confessions had the cadence of 1971, and no independent inquiry has ever
been possible. What is certain is that the affair removed the only figure in
Cuba with a military reputation comparable to the Castros', at the exact moment
the Soviet bloc was dissolving.

In April 1989 Mikhail Gorbachev came to Havana. The visit was correct and cold.
Cuba had already banned two Soviet magazines, \emph{Moscow News} and
\emph{Sputnik}, for spreading perestroika; Castro had already made the speech
that gave the period its slogan, \emph{socialismo o muerte}, socialism or
death. The Cuban position was that reform in the Soviet Union was a Soviet
matter and that Cuba would not follow. It did not follow. Nineteen months after
Gorbachev's visit the trade that had underwritten the country for eighteen
years began to stop arriving.

\begin{ledger}{abroad, 1963--1990}
\emph{Achieved.} A decisive military contribution to the defeat of South
African expansion in southern Africa, to Namibian independence, and to the
pressure that ended apartheid --- acknowledged by Mandela, paid for in Cuban
lives, and not paid for by anyone. Medical brigades in dozens of countries from
1963, the seed of Cuba's most durable asset and its eventual largest export.
A standing in the non-aligned world out of all proportion to the country's
size, which functioned as a real component of its security.

\emph{Cost.} A decade of exported insurgency that failed everywhere it was
tried and killed the people sent to try it. The endorsement of the invasion of
Czechoslovakia, made in full knowledge of what it was, as the price of oil. A
long deployment in support of the Derg through the Red Terror and the famine.
Sixteen years of war fought by conscripts whose families were told little and
whose dead came home in one ceremony. And at home, in 1980, the deliberate
mobilisation of ordinary Cubans to assault their departing neighbours ---
the one thing in this ledger for which no external necessity can be pleaded.
\end{ledger}
